\section{Conclusion}
\label{sec:conclusion}

This paper reframes the current evidence around a verification substrate gap.
Provider-side first-divergence traces can carry recoverable keyed evidence
under reviewed event-access conditions: the tested Qwen and Llama diagnostics
recover protected codewords while rejecting controls.  The same evidence does
not become a deterministic public final-text codeword in the current artifacts.
Public final-text predicates show shallow separability, but they recover zero
codeword blocks and can be searched, edited, learned, and optimized by
source-mismatch examples.

The resulting claim is intentionally narrower than an ownership-verification
claim.  The work supports cooperative trace-bound diagnostics and a negative
public final-text substrate diagnosis.  It does not support public text-only
verification success, natural evidence success, phrase-decoder success,
cryptographic provenance, or an ownership proof.  The main lesson is that
evidence strength depends on where the evidence lives.  Stronger substrates
can make verification more deterministic, but they may not travel with ordinary
natural LLM outputs; final text travels easily, but current public predicates
do not provide reliable deterministic verification.
